% design_strategy.tex — Design strategy: philosophy, methodology, and parameter derivation
% Shows the overall design philosophy, simulation-first approach, progressive testing
% rationale, team organisation, and the logical dependency chain for operating parameters.
%
% STATUS: DEFINITIVE for the step-by-step parameter derivation chain.
%         References results from path_optimization_definitive.tex (216-config sweep,
%         Pareto table) and optimization_master.tex (composite score).
%         Some figures appear in optimization_master.tex with '-app' suffix.

\section{Design Strategy}\label{sec:design-strategy}

This appendix presents the strategic rationale that shaped the project from inception through field deployment: why a simulation-first development model was adopted, how the progressive testing framework mitigated risk with limited hardware access, how the team adapted when conditions changed, and how every operating parameter was derived from evidence rather than assumption.

% =====================================================================
\subsection{Design Philosophy: Simulation-First Development}\label{sec:sim-first-philosophy}

Three interlocking principles governed the development strategy from the outset.
They were not adopted reactively---each was a deliberate engineering choice made before any code was written, based on the project's defining constraint: a student team with a single shared drone, limited flight-test opportunities, and a fixed deadline.

\paragraph{Principle 1: Prove it in software before touching hardware.}
Every algorithmic component---the 20-state finite state machine, the boustrophedon path planner, the dual-backend detection pipeline, the GPS estimation chain, and the five-layer geofence---was developed and end-to-end tested in Software-In-The-Loop (SITL) simulation on a laptop before any hardware was assembled.
SITL provides three properties that field testing cannot: repeatability (identical initial conditions), instant restart (no battery swap, no re-arming), and zero crash risk.
A simulated camera extracted FOV-correct sub-images from a georeferenced satellite map of the Fenswood site, enabling the full detect$\to$estimate$\to$land pipeline to execute against realistic terrain imagery without leaving the lab.

This approach was chosen over two alternatives.
A \emph{hardware-first} strategy---assembling the drone and iterating through flight tests---would have consumed the limited flight-day budget on debugging basic logic errors (wrong state transitions, incorrect coordinate transforms, sign inversions) that cost minutes to fix in simulation but hours in the field.
A \emph{model-based design} approach (e.g.\ Simulink/Modelica plant models with auto-generated code) would have offered higher physical fidelity but required 40--60~hours of plant modelling before the first state machine test, time that was instead spent on mission logic and detection pipeline development.

\paragraph{Principle 2: One codebase, every platform.}
A single codebase runs unmodified on the Windows development laptop, WSL/Linux, and the Raspberry~Pi~5.
Platform differences are confined to two files: \texttt{config.py} (auto-detection of serial ports, camera backends, and display availability) and \texttt{vision.py} (dual inference backend: Ultralytics on laptop, TFLite/NCNN on Pi).
No \texttt{if platform ==} guards appear in mission logic.
This eliminates an entire class of ``works on my machine'' integration failures and ensures that every simulation run exercises the same control flow that will execute in flight.

\paragraph{Principle 3: Isolate one variable per test step.}
Each test in the progressive framework introduces exactly one new variable: a new sensor, a new communication link, a new flight mode.
A failure at step~$N$ is therefore attributable to the single change introduced at step~$N$, not to an unknown interaction among multiple untested components.
This principle---borrowed from experimental design practice---transforms debugging from a combinatorial search into a linear walk through the integration ladder.

% =====================================================================
\subsection{Why This Approach Over Alternatives}\label{sec:strategy-alternatives}

The simulation-first methodology was not the only viable development strategy.
Table~\ref{tab:dev-strategy-comparison} compares three candidate approaches against the project's constraints.

\begin{table}[ht]
  \centering\compactTable
  \caption{Development strategy comparison. Scores are qualitative assessments on a 1--5 scale (5 = best fit for project constraints).}\label{tab:dev-strategy-comparison}
  \begin{tabularx}{\linewidth}{@{}X c c c@{}}
    \toprule
    \textbf{Criterion} & \textbf{Simulation-first} & \textbf{Hardware-first} & \textbf{Model-based} \\
    \midrule
    Time to first integration test      & 5 (days)   & 2 (weeks)  & 1 (months) \\
    Hardware access requirement          & 5 (none)   & 1 (daily)  & 3 (late)   \\
    Defect discovery cost                & 5 (low)    & 2 (high)   & 4 (medium) \\
    Fidelity of pre-flight validation    & 3 (medium) & 5 (high)   & 4 (high)   \\
    Team parallelism (5 members)         & 5 (full)   & 2 (serial) & 3 (partial)\\
    Risk of logic errors surviving to flight & 5 (low) & 2 (high)  & 4 (low)    \\
    \midrule
    \textbf{Total}                       & \textbf{28} & 14         & 19         \\
    \bottomrule
  \end{tabularx}
\end{table}

The simulation-first approach dominates on every criterion except pre-flight fidelity, where hardware-first testing provides direct validation of real-world effects (wind, GPS noise, motor vibration) that simulation cannot replicate.
This fidelity gap is addressed explicitly by the progressive testing ladder (Section~\ref{sec:progressive-ladder-strategy}), which introduces hardware incrementally after the software has been proven in simulation.

The hardware-first approach scores poorly on team parallelism because it serialises development around a single shared airframe: only one person can fly the drone at a time, and a crash grounds the entire team.
The simulation-first approach allows all five team members to run the full software stack simultaneously on their laptops, with integration points defined by narrow interfaces (\texttt{detect\_in\_image()} returns \texttt{(found, x, y, conf)}; \texttt{generate\_lawnmower()} returns a waypoint list) rather than by physical hardware availability.

% =====================================================================
\subsection{Progressive Testing as Risk Management}\label{sec:progressive-ladder-strategy}

The five-tier progressive testing framework (described operationally in Section~\ref{sec:testing}) is not merely a testing plan---it is the project's primary risk management instrument.
Each tier gates advancement to the next, and each failure is resolved at the cheapest possible tier before proceeding.

The cost asymmetry that motivates this structure is stark.
A sign-convention error in the geofence repulsive force (Section~\ref{sec:decisions-changed}, correction~5) cost 15~minutes to diagnose and fix in Tier~1 simulation.
Had the same error survived to Tier~5 autonomous flight, the drone would have been driven \emph{into} the SSSI no-fly zone---a Wildlife and Countryside Act 1981 violation, potential hardware loss, and a consumed flight-day slot.

Of the 15 defects discovered during the project, 12 (80\%) were caught at Tier~1 (simulation) or Tier~3 (bench test), at a combined cost of approximately 8~person-hours.
The remaining 3---all related to real-world sensor behaviour (BGR colour inversion, lens distortion magnitude, GPS timing lag)---were caught at Tier~3 and Tier~4, where the cost was a single field session rather than a failed autonomous mission.
No defects were first discovered at Tier~5.

This distribution validates the investment in simulation infrastructure: the 40+~hours spent building the SITL simulation, interactive simulator, and video replay tools were repaid many times over by avoiding field-day debugging. Figure~\ref{fig:development_timeline} visualises the full project timeline, showing how each development phase and defect discovery maps onto the progressive testing tiers.

\begin{figure}[H]
\centering
\includegraphics[width=\linewidth]{development_timeline}
\caption{Project development timeline showing the progression from simulation-first development through hardware integration. Key milestones, defect discoveries, and the weather-cancelled flight day are annotated. The majority of software development (green) was completed before hardware assembly (orange), validating the simulation-first strategy.}
\label{fig:development_timeline}
\end{figure}

% =====================================================================
\subsection{Adapting When Hardware Was Unavailable}\label{sec:adaptation-strategy}

The project's defining constraint adaptation occurred on 2026-03-11, when the scheduled first flight test was cancelled due to sustained winds exceeding the EDU-450's \SI{10}{\metre\per\second} operational limit (measured \SI{12}{\metre\per\second}, gusts to \SI{18}{\metre\per\second}).
Rather than losing the field session entirely, the team executed a pre-planned bench-testing fallback---a contingency that had been designed into the test schedule precisely because weather cancellations are a known risk for outdoor robotics projects in the UK.

The bench session yielded four results that would have been impossible to obtain without the assembled hardware:
\begin{enumerate}[nosep]
  \item \textbf{FOV calibration correction:} the datasheet focal length (\SI{7.0}{\milli\metre}) was found to be 22\% too long via a tape-measure calibration (\SI{92}{\centi\metre} visible at \SI{1}{\metre} height), yielding a corrected value of \SI{5.46}{\milli\metre}. Without this correction, GPS estimation errors at \SI{30}{\metre} would have exceeded the \SI{10}{\metre} R07 landing requirement.
  \item \textbf{Lens distortion calibration:} a 14$\times$9 checkerboard calibration produced an RMS residual of 0.399, enabling sub-pixel lens correction at \SI{1.5}{\milli\second} per frame.
  \item \textbf{Inference benchmark:} TFLite on the Pi~5 achieved \SI{206.5}{\milli\second} per frame (4.8~FPS), confirming that the search speed of \SI{8}{\metre\per\second} provides 14~frames of target visibility per flyover---well above the 3-frame minimum.
  \item \textbf{End-to-end connectivity verification:} Pi$\to$MAVProxy$\to$Cube$\to$Mission Planner chain confirmed operational, validating the MAVProxy UDP bridge architecture.
\end{enumerate}

The broader adaptation strategy extended beyond a single weather event.
Because hardware access was intermittent throughout the project (shared lab, booking conflicts, component delivery delays), the simulation-first philosophy ensured that software development was never blocked by hardware availability.
The DJI video replay tool (Section~\ref{sec:dji-proxy}) was built specifically to bridge the gap between simulation and real flight: by replaying genuine aerial footage through the identical detection pipeline, the team could validate detection performance under real lighting, texture, and altitude conditions without requiring drone access.

% =====================================================================
\subsection{Bridging the Simulation-to-Reality Gap}\label{sec:sim-real-gap-strategy}

Simulation-first development carries an inherent risk: the simulation may not be faithful enough, causing validated software to fail on real hardware.
Four specific sim-to-real gaps were identified and addressed:

\begin{enumerate}
  \item \textbf{GPS noise.}
    SITL provides near-perfect GPS with sub-metre accuracy.
    Real consumer GNSS receivers exhibit 2--3\,m CEP and 100--200\,ms timing lag.
    \emph{Mitigation:} the interactive simulator includes configurable Ornstein--Uhlenbeck GPS drift (up to \SI{5}{\metre} $\sigma$) and the Kalman filter estimator was validated against these injected errors. DJI video replay confirmed a CEP50 of \SI{2.3}{\metre} under real conditions.

  \item \textbf{Detection under real conditions.}
    Synthetic training data and simulated camera views cannot reproduce real-world lighting variation, partial occlusion, or terrain texture.
    \emph{Mitigation:} the model was retrained on a mixed dataset (300~synthetic + 16~real aerial photos + 50~negatives) and validated via DJI video replay, which confirmed 87\% detection rate on frames where the target exceeded $20 \times 20$\,px.

  \item \textbf{Motion blur and vibration.}
    By default, SITL camera views are sharp synthetic crops; real aerial imagery exhibits motion blur proportional to speed and altitude.
    \emph{Mitigation:} three CLI flags inject degradation into the SITL camera view: \texttt{-{}-blur} applies physics-based directional motion blur (kernel oriented along the yaw angle, sized by translational ground speed and rotational yaw rate at a simulated 5\,ms exposure); \texttt{-{}-shake} adds per-frame random pixel translation modelling vibration; and \texttt{-{}-noise} injects Gaussian noise into GPS (${\sim}2.8$\,m std), altitude (${\sim}0.5$\,m std), and attitude (${\sim}0.5^{\circ}$ std). The interactive simulator provides equivalent camera shake and GPS drift controls. Additionally, the IMX296 global shutter eliminates rolling-shutter distortion, and the speed schedule (\SI{8}{\metre\per\second} at \SI{35}{\metre}) limits inter-frame ground displacement to \SI{1.67}{\metre}, well within the \SI{32}{\metre} footprint width.

  \item \textbf{Wind disturbance.}
    No simulation level models wind effects on flight stability or lane tracking accuracy.
    \emph{Mitigation:} the 20\% swath overlap provides \SI{6.4}{\metre} of lateral margin against the \SI{4.6}{\metre} RSS displacement budget (GPS error + crosswind + roll coupling), and the NFZ waypoint buffer (\SI{30}{\metre}) provides a $2.2\times$ safety factor above the \SI{23}{\metre} RSS requirement.
\end{enumerate}

Each gap was treated as a quantified risk with a specific mitigation, rather than an unacknowledged limitation.
The mitigations are conservative by design: if real conditions prove worse than assumed, the system degrades gracefully (slower search, wider margins, more conservative detection threshold) rather than failing catastrophically.

% =====================================================================
\subsection{Team Organisation and Workstream Structure}\label{sec:team-strategy}

The five-member team was organised into parallel workstreams to maximise development velocity while maintaining integration discipline:

\begin{itemize}[nosep]
  \item \textbf{Software and CV} (Dmytro): autonomous software stack, detection pipeline, simulation environment, training data pipeline, and progressive testing framework.
  \item \textbf{Flight Dynamics and Testing} (Robin): state machine testing, flight parameter tuning, simulation validation, and Cube flight-mode configuration.
  \item \textbf{RC and Pilot Operations} (Teammate~3): transmitter configuration, kill-switch setup, manual override procedures, and designated pilot for all flight tests.
  \item \textbf{Hardware Integration} (Teammate~4): airframe assembly, wiring, camera and GPS mounting, compass calibration, and payload CG verification.
  \item \textbf{Project Management} (Teammate~5): scheduling, flight-slot booking, risk assessments, report compilation, and deliverable tracking.
\end{itemize}

This structure was chosen over two alternatives.
A \emph{full-stack} model (every member works on everything) would have created merge conflicts and knowledge fragmentation across a 10,000+~line codebase.
A \emph{pair-programming} model would have doubled the person-hours on each task without proportional quality improvement, given the team's heterogeneous skill profiles (software engineering, aerodynamics, mechanical assembly, project management).

The workstream model's key risk---integration failures at stream boundaries---was mitigated by narrow, well-defined interfaces.
The CV stream delivers \texttt{(found, x, y, conf)}.
The path planner delivers a waypoint list.
The hardware stream delivers a working serial link.
Integration testing (bench tests, passive flights) was conducted jointly, ensuring that no subsystem was developed in isolation from the system it would connect to.

% =====================================================================
\subsection{Parameter Derivation Chain}\label{sec:param-derivation-chain}

The preceding strategic choices---simulation-first development, progressive testing, single-codebase architecture---created the infrastructure within which every operating parameter could be derived from evidence rather than assumed.
The operating parameters in Table~\ref{tab:pareto-final} were not selected independently.
Each decision constrains the next, forming a directed dependency chain from the physical target through to the final mission plan.
The remainder of this appendix traces that chain step by step.

% =====================================================================
\subsection{Step 1: Maximum Detection Altitude}\label{sec:step-max-alt}

The reasoning begins with the target.
The mannequin is \SI{1.8}{\metre} tall and \SI{0.5}{\metre} wide.
At altitude~$h$, the target projects onto the camera sensor with a height in pixels of:

\begin{equation}\label{eq:target-px}
  \mathrm{px}_{\mathrm{target}} = \frac{d_{\mathrm{target}} \cdot f_{\mathrm{px}}}{h}
\end{equation}

\noindent where $f_{\mathrm{px}} = f / w_s \times W_{\mathrm{px}} = 5.46/5.02 \times 1456 = 1584$\,pixels is the focal length in pixel units, and $d_{\mathrm{target}}$ is the target dimension.
YOLOv8n resizes the $1456 \times 1088$ sensor frame to $640 \times 640$; the uniform scale factor is $640/\max(1456, 1088) = 640/1456 = 0.4396$:

\begin{equation}\label{eq:model-px}
  \mathrm{px}_{\mathrm{model}} = \frac{d_{\mathrm{target}} \cdot f_{\mathrm{px}}}{h} \times \frac{640}{1456}
\end{equation}

Table~\ref{tab:px-vs-alt} traces this function across the altitude range.

\begin{table}[ht]
  \centering\compactTable
  \caption{Target pixel extent versus altitude. Sensor column = pixels in the native $1456\!\times\!1088$ frame; model-input columns = after resize to $640\!\times\!640$ (scale $= 640/1456 = 0.4396$). Below \SI{20}{px} height in the model input, detection becomes unreliable.}\label{tab:px-vs-alt}
  \begin{tabularx}{\linewidth}{@{}c c c c X@{}}
    \toprule
    \textbf{Altitude (m)} & \textbf{Sensor (px)} & \textbf{Model input (px)} & \textbf{Width model (px)} & \textbf{Assessment} \\
    \midrule
    20 & 142 & 63 & 18 & Excellent --- large, high confidence \\
    35 & 81  & 36 & 10 & Comfortable --- clear margin above minimum \\
    50 & 57  & 25 & 7  & Marginal width --- feature extraction degrades \\
    60 & 48  & 21 & 6  & Below width comfort threshold \\
    85 & 34  & 15 & 4  & At absolute minimum; unreliable \\
    \bottomrule
  \end{tabularx}
\end{table}

Two thresholds emerge.
The \emph{absolute} detection limit occurs when the target height drops below approximately \SI{20}{px} in the model input ($h \approx \SI{70}{\metre}$), at which point the feature map representation is too coarse for reliable bounding box regression.
The \emph{practical} limit is set by target \emph{width}: at \SI{50}{\metre} the mannequin's width is only 7~model pixels, below the 10-pixel comfort threshold for width-dependent features (limb visibility, aspect ratio cues).
Empirically, detection confidence from the retrained \texttt{sar\_v2\_1088} model drops below 0.5 at approximately \SI{60}{\metre}.

Regulatory constraint R03 limits altitude to \SI{50}{\metre}.
Selecting $h = \SI{35}{\metre}$ provides a $\times$1.4 margin on target width (10\,px vs 7\,px at \SI{50}{\metre}) and a $\times$1.8 margin on target height (36\,px vs 20\,px minimum), creating headroom for real-world degradation modes: partial occlusion, unusual pose, non-ideal lighting, and altitude excursions during manoeuvres.

\textbf{Decision:} $h = \SI{35}{\metre}$.
This sacrifices the energy-optimal \SI{50}{\metre} (\SI{8.2}{Wh}) for a comfortable detection margin (\SI{12.6}{Wh}---still only \SI{10.9}{\percent} of battery capacity).

% =====================================================================
\subsection{Step 2: Maximum Speed at Chosen Altitude}\label{sec:step-max-speed}

With altitude fixed, the camera's along-track ground footprint determines how long the target remains in view.
At \SI{35}{\metre}:

\begin{equation}\label{eq:footprint-along}
  H_g = \frac{w_s \cdot h}{f} \times \frac{H_{\mathrm{px}}}{W_{\mathrm{px}}} = \frac{5.02 \times 35}{5.46} \times \frac{1088}{1456} = \SI{24.1}{\metre}
\end{equation}

At the TFLite inference rate of $f_{\mathrm{inf}} = \SI{4.8}{FPS}$ (benchmarked on Pi~5, Section~\ref{sec:cv:inference:pipeline}), the drone advances $\Delta x = v / f_{\mathrm{inf}}$ metres between frames.
For reliable detection, the target must appear in at least $n_{\min}$ independent inference frames.
The maximum speed is therefore:

\begin{equation}\label{eq:max-speed}
  v_{\max} = \frac{H_g}{n_{\min} / f_{\mathrm{inf}}} = \frac{H_g \cdot f_{\mathrm{inf}}}{n_{\min}}
\end{equation}

If we require $n_{\min} = 3$ frames (one detection above threshold triggers investigation), then $v_{\max} = 24.1 \times 4.8 / 3 = \SI{38.6}{\metre\per\second}$---far above drone capability.
However, at high frame rates consecutive frames are highly correlated: at \SI{8}{\metre\per\second} the inter-frame advance is \SI{1.67}{\metre}, yielding $91\%$ spatial overlap between adjacent frames.
These are not independent observations.

A conservative requirement of $n_{\min} = 10$ raw frames ensures that at least 3--4 contain substantially different views of the target (different pixel positions, different background context):

\begin{equation}
  v_{\max,\mathrm{conservative}} = \frac{24.1 \times 4.8}{10} = \SI{11.5}{\metre\per\second}
\end{equation}

The altitude-dependent speed schedule (Section~\ref{sec:moo}) assigns $v(\SI{35}{\metre}) = \SI{8.0}{\metre\per\second}$, which yields:

\begin{equation}
  n_f = \frac{H_g}{v} \times f_{\mathrm{inf}} = \frac{24.1}{8.0} \times 4.8 = 14.5 \approx 14 \text{ frames}
\end{equation}

\noindent This provides a $\times$1.4 margin above the conservative 10-frame requirement and a $\times$4.7 margin above the absolute 3-frame minimum.
The cumulative detection probability over 14 frames, assuming per-frame probability $p = 0.95$ (conservative for 48-pixel targets), is $P_d = 1 - (1-0.95)^{14} > 0.9999$.

\textbf{Decision:} $v = \SI{8}{\metre\per\second}$ at \SI{35}{\metre}, following the altitude-dependent linear schedule.

% =====================================================================
\subsection{Step 3: Minimise Energy via Scan Angle}\label{sec:step-scan-angle}

With $h$ and $v$ fixed, the remaining free variable for the boustrophedon pattern is the scan angle~$\theta$---the orientation of parallel scan lines relative to north.
This controls the number of scan lines (and therefore U-turns) required to cover the polygon.

The 216-configuration parametric sweep (Section~\ref{sec:sweep}) tested 36 angles at each of 6 altitudes.
At $h = \SI{35}{\metre}$, the results are unambiguous: scan angle $\theta = \SI{70}{\degree}$ (alignment with the polygon's longest edge, the northeast--southwest axis) minimises the number of scan lines from 10 (at \SI{0}{\degree}, north--south) to 6, and reduces U-turns from 9 to 5.

The energy savings are substantial.
At the optimal angle: $E = \SI{12.6}{Wh}$.
At the worst angle (\SI{95}{\degree}, perpendicular to the long axis): $E = \SI{25.4}{Wh}$---a $2\times$ penalty for the same coverage.
The energy valley at $\theta \approx \SI{70}{\degree}$ is visible in Figure~\ref{fig:energy-heatmap} across all altitudes, confirming that scan angle alignment is geometry-dependent and altitude-invariant.

This angle also has a fortunate geometric property: the scan lines run roughly \emph{perpendicular} to the SSSI boundary on the northeast edge.
This means each scan line clips the NFZ slow zone only at its endpoint, rather than running parallel to the boundary (which would force the entire line to be traversed at reduced speed).
The NFZ energy penalty at \SI{70}{\degree} is approximately \SI{15}{\percent}; at \SI{0}{\degree} (north--south, running along the SSSI edge), the penalty exceeds \SI{40}{\percent}.

\textbf{Decision:} $\theta = \SI{70}{\degree}$, aligned to the polygon's longest edge via \texttt{cv2.minAreaRect}.

% =====================================================================
\subsection{Step 4: NFZ Safety Margins}\label{sec:step-nfz-margins}

The SSSI no-fly zone on the polygon's northeast edge requires a buffer that keeps the entire camera footprint---not just the drone's position---outside the protected area.
At \SI{35}{\metre} altitude, the cross-track ground footprint is:

\begin{equation}\label{eq:footprint-cross}
  W_g = \frac{w_s \cdot h}{f} = \frac{5.02 \times 35}{5.46} = \SI{32.2}{\metre}
\end{equation}

The camera images a strip extending $W_g / 2 = \SI{16.1}{\metre}$ either side of nadir.
To guarantee the camera footprint does not overlap the SSSI, the drone must maintain at least \SI{16.1}{\metre} lateral distance from the boundary.
Three additional error sources widen this margin:

\begin{enumerate}
  \item \textbf{GPS position error}: CEP$_{95}$ = \SI{3}{\metre} for consumer-grade GNSS.
  \item \textbf{Reaction time}: at $v = \SI{8}{\metre\per\second}$, a \SI{2}{\second} reaction window covers \SI{16}{\metre}.
  \item \textbf{Wind gust displacement}: estimated \SI{2}{\metre} lateral excursion in \SI{15}{kt} crosswind.
\end{enumerate}

The root-sum-square (RSS) of these independent errors is:

\begin{equation}
  d_{\mathrm{margin}} = \sqrt{16.1^2 + 3^2 + 16^2 + 2^2} = \SI{23.1}{\metre}
\end{equation}

The system implements a five-layer defence (see Section~\ref{sec:geofencing} for full details):

\begin{itemize}
  \item \textbf{\SI{30}{\metre} waypoint buffer} (\texttt{NFZ\_WAYPOINT\_BUFFER\_M}): any waypoint within \SI{30}{\metre} of the SSSI boundary is excluded at plan time.
  \item \textbf{\SI{20}{\metre} speed ramp} (\texttt{NFZ\_SLOW\_ZONE\_M}): linear deceleration from \SI{3}{\metre\per\second} to \SI{0.3}{\metre\per\second} within this zone.
  \item \textbf{Repulsive vector field}: constant-magnitude push away from the boundary when within \SI{3}{\metre}.
  \item \textbf{\SI{3}{\metre} hard boundary} (\texttt{NFZ\_HARD\_BOUNDARY\_M}): automatic switch to MANUAL mode.
  \item \textbf{Firmware geofence}: independent Cube Orange autopilot geofence triggers LOITER/RTL as a final backstop.
\end{itemize}

The combined margin of \SI{30}{\metre} (waypoint buffer) + \SI{20}{\metre} (speed ramp) = \SI{50}{\metre} exceeds the RSS requirement of \SI{23}{\metre} by a factor of $2.2\times$.
As a cross-check, the one-third rule (buffer $\geq W_g / 3$) requires \SI{10.7}{\metre}---our \SI{30}{\metre} buffer exceeds this by $2.8\times$.

The cost of this conservatism is a \SI{4}{\percent} coverage gap in the strip adjacent to the SSSI.
This strip lies within a fenced wildlife reserve with no public access, making the probability of a casualty there negligibly small.

\textbf{Decision:} \SI{30}{\metre} waypoint buffer, \SI{20}{\metre} speed ramp, \SI{3}{\metre} hard cut-off.
Accept \SI{96}{\percent} polygon coverage.

% =====================================================================
\subsection{Step 5: Swath Overlap}\label{sec:step-overlap}

Adjacent scan lines must overlap sufficiently to prevent coverage gaps caused by lateral position error.
Three sources of lane misalignment are relevant:

\begin{enumerate}
  \item \textbf{GPS cross-track drift}: $\sigma_{\mathrm{GPS}} = \SI{3}{\metre}$ (consumer GNSS CEP).
  \item \textbf{Crosswind displacement}: $\delta_{\mathrm{wind}} \approx \SI{1.5}{\metre}$ lateral shift in \SI{10}{kt} crosswind at \SI{8}{\metre\per\second}.
  \item \textbf{Attitude roll coupling}: a \SI{5}{\degree} roll displaces the footprint centre by $h \tan(5\degree) = 35 \times 0.087 = \SI{3.1}{\metre}$.
\end{enumerate}

Assuming independence, the worst-case combined displacement is:

\begin{equation}\label{eq:lane-error}
  \delta_{\mathrm{total}} = \sqrt{3^2 + 1.5^2 + 3.1^2} = \SI{4.6}{\metre}
\end{equation}

At \SI{35}{\metre} altitude with \SI{20}{\percent} overlap, the overlap margin in metres is:

\begin{equation}
  m_{\mathrm{overlap}} = 0.20 \times W_g = 0.20 \times 32.2 = \SI{6.4}{\metre}
\end{equation}

This covers the worst-case displacement with \SI{1.8}{\metre} to spare ($6.4 - 4.6 = \SI{1.8}{\metre}$).
Increasing to \SI{30}{\percent} overlap would add an extra scan line (7 instead of 6), increasing flight time by approximately \SI{10}{\percent} for only \SI{3.2}{\metre} of additional margin---beyond the point of diminishing returns given the RSS error model.

\textbf{Decision:} \SI{20}{\percent} swath overlap.

% =====================================================================
\subsection{Step 6: Diagonal vs Aligned Scanning}\label{sec:step-diagonal}

A diagonal scan orientation (e.g., \SI{45}{\degree} to the wind) has been suggested in the coverage planning literature as a way to reduce crosswind-induced lane drift, since the wind component perpendicular to the scan direction is reduced~\cite{galceran2013survey}.
For the AENGM0074 polygon, this would mean rotating the scan lines approximately \SI{25}{\degree} from the longest-edge alignment.

The cost is significant: rotating from \SI{70}{\degree} to \SI{45}{\degree} increases the number of scan lines from 6 to approximately 8 (a \SI{33}{\percent} increase), adding 2 U-turns and approximately \SI{300}{\metre} of additional path length.
The resulting energy cost rises from \SI{12.6}{Wh} to approximately \SI{16.4}{Wh}---a \SI{30}{\percent} penalty.

The benefit is marginal: the crosswind displacement (\SI{1.5}{\metre}) is already covered by the \SI{6.4}{\metre} overlap margin with \SI{4.9}{\metre} to spare.
Even doubling the crosswind to \SI{20}{kt} (\SI{3}{\metre} displacement) leaves \SI{3.4}{\metre} of margin within the existing \SI{20}{\percent} overlap.

\textbf{Decision:} maintain longest-edge alignment (\SI{70}{\degree}).
The \SI{20}{\percent} overlap already absorbs crosswind effects without the \SI{30}{\percent} energy penalty of diagonal scanning.

% =====================================================================
\subsection{The Dependency Chain}\label{sec:dependency-chain}

The six steps above form a directed acyclic graph of parameter dependencies, where each decision is conditioned on the outputs of previous steps:

\begin{enumerate}
  \item \textbf{Target size} $\rightarrow$ \textbf{altitude}:
    the mannequin's \SI{1.8}{\metre} height and the 20-pixel model-input minimum determine $h_{\max} = \SI{70}{\metre}$;
    the 10-pixel width comfort threshold sets $h_{\mathrm{comfort}} = \SI{50}{\metre}$;
    a \SI{30}{\percent} safety margin yields $h = \SI{35}{\metre}$.

  \item \textbf{Altitude} $\rightarrow$ \textbf{footprint} $\rightarrow$ \textbf{speed}:
    $h = \SI{35}{\metre}$ gives a \SI{24}{\metre} along-track footprint;
    at \SI{4.8}{FPS}, requiring 10+ raw frames caps speed at \SI{11.5}{\metre\per\second};
    the altitude-dependent schedule assigns $v = \SI{8}{\metre\per\second}$.

  \item \textbf{Altitude + speed} $\rightarrow$ \textbf{scan angle}:
    with $h$ and $v$ fixed, the 216-configuration sweep identifies $\theta = \SI{70}{\degree}$ as the energy minimum (longest-edge alignment, fewest U-turns).

  \item \textbf{Altitude} $\rightarrow$ \textbf{footprint width} $\rightarrow$ \textbf{NFZ margin}:
    the \SI{32}{\metre} footprint at \SI{35}{\metre} requires at least \SI{16}{\metre} standoff;
    GPS error, reaction time, and wind push this to \SI{23}{\metre} RSS;
    the \SI{30}{\metre} buffer provides a $2.2\times$ safety factor.

  \item \textbf{Footprint + error model} $\rightarrow$ \textbf{overlap}:
    \SI{4.6}{\metre} RSS lane error is covered by \SI{6.4}{\metre} (\SI{20}{\percent}) overlap margin.

  \item \textbf{Energy budget} validates the chain:
    \SI{12.6}{Wh} is \SI{10.9}{\percent} of the \SI{115.4}{Wh} battery, leaving \SI{14}{\minute} of flight endurance---sufficient for the \SI{1.9}{\minute} search plus transit, descent, verification, and landing.
\end{enumerate}

\begin{figure}[ht]
  \centering
  \includegraphics[width=\linewidth]{optimization_network}
  \caption{Optimisation variable dependency network. Decision variables (blue) cascade through derived quantities (white) to performance metrics (green/red). Arrow thickness indicates influence strength.}
  \label{fig:optimization-network}
\end{figure}

\begin{figure}[ht]
  \centering
  \includegraphics[width=0.85\linewidth]{sensitivity_matrix}
  \caption{Sensitivity matrix showing the influence of each input parameter on every derived quantity. Darker cells indicate stronger coupling, revealing which measurements most critically affect mission performance.}
  \label{fig:sensitivity-matrix}
\end{figure}

\begin{figure}[ht]
  \centering
  \includegraphics[width=0.85\linewidth]{sensitivity_spider}
  \caption{Spider plot of normalised parameter sensitivities. Each axis represents a derived performance metric; the polygon area indicates the overall sensitivity footprint of the design point.}
  \label{fig:sensitivity-spider}
\end{figure}

\FloatBarrier

\begin{figure}[ht]
  \centering
  \includegraphics[width=0.85\linewidth]{coupling_matrix}
  \caption{Variable coupling matrix showing pairwise interaction effects. The altitude--speed pair exhibits the only strong coupling, jointly determining coverage rate and detection probability. Off-diagonal cells indicate which performance metric is affected by each variable pair.}
  \label{fig:coupling-matrix}
\end{figure}

\begin{figure}[ht]
  \centering
  \includegraphics[width=0.85\linewidth]{altitude_speed_tradeoff}
  \caption{Altitude--speed trade-off space. Colour shows composite mission score (40\% detection $+$ 35\% coverage $+$ 25\% energy). Blue contours: detection probability. Orange contours: coverage time. The selected operating point (35\,m, 8\,m/s) sits above the 95\% detection contour in the high-score region.}
  \label{fig:alt-speed-tradeoff}
\end{figure}

\begin{figure}[ht]
  \centering
  \includegraphics[width=0.88\linewidth]{function_chain}
  \caption{Function chain diagram tracing the flow from physical inputs (target size, sensor parameters, inference rate) through intermediate quantities (GSD, pixel extent, dwell time) to the five mission performance objectives. Each arrow represents a functional dependency with the governing equation labelled.}
  \label{fig:function-chain}
\end{figure}

\FloatBarrier

\begin{figure}[ht]
  \centering
  \includegraphics[width=0.88\linewidth]{tornado_sensitivity}
  \caption{Tornado diagram showing the sensitivity of the composite mission score to $\pm$20\% perturbations in each input parameter. Altitude and inference FPS have the largest influence, confirming that the detection--energy trade-off dominates. Parameters are sorted by total bar length (influence magnitude).}
  \label{fig:tornado-sensitivity}
\end{figure}

\noindent Every parameter is traceable to a physical measurement (target dimensions, sensor specifications, inference benchmarks) or a regulatory constraint (altitude ceiling, NFZ boundaries). Figures~\ref{fig:sensitivity-matrix} and~\ref{fig:sensitivity-spider} quantify single-variable sensitivities, Figure~\ref{fig:coupling-matrix} reveals pairwise interactions, and Figure~\ref{fig:alt-speed-tradeoff} maps the dominant altitude--speed trade-off space. Figure~\ref{fig:optimization-network} visualises the full dependency chain.
No parameter was assumed or rounded for convenience.
The chain can be re-evaluated instantly if any input changes---for example, a faster inference backend would relax the speed constraint, and a different target size would shift the altitude ceiling---making the design robust to future system upgrades.
